%==============================================================================
%  Lautaro Chittaro — Curriculum Vitae
%  Title-based layout (centered header + ruled section headings)
%==============================================================================
\documentclass[12pt]{article}

\usepackage[T1]{fontenc}
\usepackage[utf8]{inputenc}
\usepackage[letterpaper,top=1in,bottom=1in,left=1in,right=1in]{geometry}
\usepackage{newtxtext}           % Times-like serif for the body
\usepackage{newtxmath}
\usepackage{xcolor}
\usepackage{ragged2e}
\usepackage{tabularx}
\usepackage[hidelinks]{hyperref}

% Minimal \needspace: force a page break if less than #1 vertical space remains,
% so a heading/title is never stranded at the foot of a page (needspace.sty is
% not installed in this TeX tree).
\makeatletter
\newcommand{\needspace}[1]{%
  \begingroup
    \setlength{\dimen@}{#1}%
    \dimen@ii=\pagegoal \advance\dimen@ii-\pagetotal
    \ifdim\dimen@>\dimen@ii \newpage \fi
  \endgroup}
\makeatother

% ---- Colours ----------------------------------------------------------------
\definecolor{bocconiblue}{HTML}{0047AB} % cobalt-blue accent for links
\definecolor{rulegray}{HTML}{9A9A9A}    % soft gray for section rules

\hypersetup{
  colorlinks=true,
  urlcolor=bocconiblue,
  linkcolor=bocconiblue,
  citecolor=bocconiblue,
  pdftitle={Lautaro Chittaro - Curriculum Vitae},
  pdfauthor={Lautaro Chittaro}
}

\setlength{\parindent}{0pt}
\linespread{1.04}

% ---- Avoid widow / orphan lines ---------------------------------------------
\widowpenalty=10000
\clubpenalty=10000
\displaywidowpenalty=10000
\raggedbottom

% ---- Unified left column ----------------------------------------------------
% Every section body is indented to \bodyindent; years live in the gutter to
% the left of it, so all body text lines up in a single vertical column while
% headings and rules stay flush with the page margin.
\newlength{\bodyindent}\setlength{\bodyindent}{1.65cm}
\newlength{\bodywidth}\setlength{\bodywidth}{\dimexpr\textwidth-\bodyindent\relax}

% ---- Ruled section heading (flush left, full width) -------------------------
\newcommand{\cvsec}[1]{%
  \par
  \vspace{1.05em}%
  {\scshape\large #1}\par
  \vspace{0.15em}%
  {\color{rulegray}\hrule height 0.3pt}%
  \vspace{0.6em}%
}

% ---- Dated entry: content flush left, year right-aligned (wraps safely) -----
\newcommand{\yr}[2]{%
  \par\noindent
  \begin{tabularx}{\textwidth}{@{}X@{\hspace{1.5em}}r@{}}
    #2 & {\small #1}
  \end{tabularx}\par}
% Year-first variant (for list-style sections such as seminars).
\newcommand{\yrleft}[2]{%
  \par\noindent
  {\small #1}\hspace{0.6em}#2\par}

% ---- Building blocks --------------------------------------------------------
% Plain (non-link) title: bold black. \needspace keeps the title from being
% stranded alone at the foot of a page (it moves down with its abstract).
\newcommand{\entrytitle}[1]{\needspace{5\baselineskip}{\bfseries #1}}
% Linked title: bold + accent colour, wraps a real hyperlink.
\newcommand{\entrytitlelink}[2]{\needspace{5\baselineskip}\href{#1}{\bfseries\color{bocconiblue} #2}}
\newcommand{\coauthors}[1]{{\normalfont #1}\par}
\newcommand{\abstractpar}[1]{\par\vspace{0.25em}{\small #1\par}}

\begin{document}
\pagestyle{empty}

% ---- Centered header --------------------------------------------------------
\begin{center}
  {\LARGE\scshape Lautaro Chittaro}\\[0.5em]
  Department of Finance, Bocconi University\\
  Via Roentgen 1, 20136, Milan, Italy\\[0.35em]
  \href{mailto:lautaro.chittaro@unibocconi.it}{lautaro.chittaro@unibocconi.it} \quad
  \href{https://lautarochittaro.github.io}{lautarochittaro.github.io}
\end{center}

% ---- Academic Positions -----------------------------------------------------
\cvsec{Academic Positions}
\yr{2026}{Assistant Professor of Finance, Bocconi University}

% ---- Education --------------------------------------------------------------
\cvsec{Education}
\yr{2026}{Ph.D. in Economics, Stanford University}
\yr{2016}{M.A. in Economics, Universidad de San Andr\'es}
\yr{2015}{B.A. in Economics, Universidad de Buenos Aires}

% ---- References (hidden) ----------------------------------------------------
% \cvsec{References}
% \noindent
% \begin{minipage}[t]{0.5\bodywidth}
%   \textbf{Monika Piazzesi} (co-primary advisor)\\
%   Dept.\ of Economics, Stanford University\\
%   \href{mailto:piazzesi@stanford.edu}{piazzesi@stanford.edu}
% \end{minipage}%
% \begin{minipage}[t]{0.5\bodywidth}
%   \textbf{Martin Schneider} (co-primary advisor)\\
%   Dept.\ of Economics, Stanford University\\
%   \href{mailto:schneidr@stanford.edu}{schneidr@stanford.edu}
% \end{minipage}
%
% \vspace{0.7em}
% \textbf{Shoshana Vasserman}\\
% Grad.\ Sc.\ of Business, Stanford University\\
% \href{mailto:svass@stanford.edu}{svass@stanford.edu}

% ---- Fields -----------------------------------------------------------------
\cvsec{Fields}
Financial Economics, Macroeconomics, Industrial Organization

% ---- Working Papers (includes the Job Market Paper) -------------------------
\cvsec{Working Papers}
\entrytitlelink{https://lautarochittaro.github.io/assets/papers/fogape.pdf}{Selection in Crisis Lending: Evidence from Chile's Government-Guaranteed Loans}\par
\coauthors{with Crist\'ian S\'anchez}
\abstractpar{We study the long-run effectiveness of government-guaranteed loan programs
implemented during recent crises. Using administrative and loan application data from the
Central Bank of Chile, we track firm defaults five years after the COVID-19 shock. Our
instrumental-variable estimates show that these loans postponed defaults for two years but did
not reduce total defaults in the long run. Banks used private information to direct credit
toward firms that would have been safer even without the program. To assess the welfare
implications of the delayed defaults and banks' selection of safer firms, we build a dynamic
model of heterogeneous entrepreneurs disciplined by our causal estimates. The program generated
welfare gains 21\% above its fiscal cost, with limited rents for banks and modest increases in
aggregate risk-taking. Younger firms are the most cost-effective group to support, yet they are
the least likely to be approved because their growth relies on leverage, increasing default
risk. A budget-neutral redesign that raises guarantees for younger firms and reduces them for
the rest could increase welfare by 6pp.}

\vspace{0.7em}
\entrytitlelink{https://office365stanford-my.sharepoint.com/:b:/g/personal/chittaro_stanford_edu/EVvnhDn6dbZGk3inqLAbdsMBBV_AK1OOnfWAP_kEyaT6FQ?e=4hMK6z}{Asset Returns as Carbon Taxes}\par
\coauthors{with Monika Piazzesi, Martin Schneider and Marcelo Sena}
\abstractpar{In frictionless financial markets, a carbon tax on energy users provides the same
incentives as a \textit{replicating return schedule} that depends on firms' emission
intensities, defined as scope 1 emissions relative to enterprise value. We use this result to
interpret pollution premia measured by recent empirical studies and conclude that markets
currently provide only modest incentives. Replicating a serious carbon tax requires high returns
in the right tail of the emission intensity distribution. With heterogeneous investors, such
returns are not sustainable unless essentially everyone perceives large nonpecuniary costs from
holding dirty capital. Substantial emission reductions can be achieved, however, when even a
small share of investors perceive nonpecuniary \textit{benefits} from owning clean electricity
capital.}

\vspace{0.7em}
\entrytitlelink{https://marcelosena.github.io/marcelosenajmp.pdf}{Pricing and Risk in Sovereign Green Debt: Evidence from Chile}\par
\coauthors{with Marcelo Sena}
\abstractpar{We study the pricing of sovereign green bonds using Chile's pioneering green bond
program and its cross-design issuance. Employing a panel of Chilean U.S.-dollar bonds, we
estimate no-arbitrage pricing kernels for green and conventional bonds. The results reveal a
declining greenium across maturities, driven by the higher interest-rate risk exposure of green
bonds. We find no evidence of investor segmentation or liquidity differences between green and
conventional bonds. Instead, we explain the observed pricing patterns through a
representative-agent asset-pricing model in which investors derive nonpecuniary benefits from the
real value of their green bond holdings. During high-inflation periods, as observed in our
sample, the real value of green bond portfolios deteriorates, making the convenience service
they provide scarcer and more valuable. This positive correlation between green convenience yield
and inflation generates a risk premium that compresses the greenium especially at longer
maturities, producing a downward-sloping greenium term structure.}

% ---- Work in Progress -------------------------------------------------------
\cvsec{Work in Progress}
\entrytitle{The Cost of Port Disruptions: Evidence from U.S. Containerized Trade}\par
\coauthors{with Stephen Redding, Janet Stefanov, and Shoshana Vasserman}
\abstractpar{How costly are disruptions to activity at U.S. ports? To answer this question, we
estimate a model of demand for importers of different types of products who choose which maritime
port (if any) to use for their container imports. Our estimator leverages a nearly comprehensive
panel of maritime imports to the U.S. between 2020 and 2025, linking customs records, granular
GPS pings from container-shipping fleets, origin-destination-level shipping prices, and a number
of additional data sources. Using variation in prices and processing times from localized
slowdowns and disruptions at different ports, our model rationalizes importers' port choices as a
function of daily origin-destination prices, travel time at sea and on land, port congestion, and
sticky product-origin-destination preferences at the time of each shipment's observed departure.
Our estimates allow us to predict the economic incidence of a short-term disruption at a subset of
ports, such as a general strike. We use this to discuss the potential of different policies to
support supply chain resilience.}

% ---- Academic Publications --------------------------------------------------
\cvsec{Pre-PhD Publications}
\entrytitlelink{https://latinaer.org/index.php/laer/article/view/38}{From Bad to Worse: the
Economic Impact of COVID-19 in Developing Countries. Evidence from Venezuela}\par
\coauthors{with Germ\'an Caruso, Mar\'ia Emilia Cucagna, Luis Pedro Espa\~na}
\textit{Latin American Economic Review}, 2021\par
\abstractpar{Policy responses to COVID-19 affected the dynamic of economic growth and labor
markets worldwide, hitting economically harder on developing countries. These policies involved
economic lockdowns that included the shutdown of the main statistical exercises, making it almost
impossible to assess the breadth and variety of their effects. Using a phone survey, this paper
examines the impact of the quarantine implemented in Venezuela on labor market outcomes. The
identification strategy exploits the exogenous variation in the severity of the lockdown in
different regions of the country. The main result indicates a 16.5 percentage points reduction in
employment, while in regions with severe lockdowns the reduction has been 13.8 p.p.\ larger. In
particular, the self-employed and informally employed were hard hit by the lockdown. To cope with
this effect, households sold their productive assets, reduced their savings, sought for
alternative income sources and looked for help from relatives. This paper does not find a
differential effect on the number of COVID-19 cases in more severe lockdown settings. Results are
robust to endogenous migration and alternative specifications.}

% ---- Conferences and Seminars -----------------------------------------------
\cvsec{Conferences and Seminars}
\yrleft{2026}{ITAM, PUC-Rio, FGV EPGE, Universidad de los Andes, University of Miami,
  Boston College, Bocconi University, CEMFI, Universidad de Chile, EEA-ESEM,
  Universidad de Buenos Aires, Universidad Torcuato Di Tella, Banco Central de Chile}
\yrleft{2025}{NBER Summer Institute, Economic Fluctuations and Growth}

% ---- Other Positions --------------------------------------------------------
\cvsec{Other Positions}
\yr{2022-25}{Department of Economics, Stanford University}
Research Assistant for Martin Schneider and Monika Piazzesi\par
\vspace{0.4em}
\yr{2022-25}{Graduate School of Business, Stanford University}
Research Assistant for Shoshana Vasserman\par
\vspace{0.4em}
\yr{2020}{World Bank, Washington DC}
Short-term Consultant\par
\vspace{0.4em}
\yr{2017-20}{Ministry of Production, Argentina}
Trade Policy Advisor\par

% ---- Teaching ---------------------------------------------------------------
\cvsec{Teaching}
Macroeconomics II (Ph.D.), Stanford University\\
Economic Forecasting (Undergrad), Stanford University\\
Macroeconomics I, II (Undergrad), Universidad de Buenos Aires\\
Industrial Organization (Undergrad), Universidad de Buenos Aires\\
National Accounts (Undergrad), Universidad de Buenos Aires

% ---- Awards & Fellowships ---------------------------------------------------
\cvsec{Awards \& Fellowships}
\yr{2024}{Gale and Steve Kohlhagen Fellowship, Stanford University}
\yr{2020-22}{The Alejandro and Lida Zaffaroni Fellowship, Stanford University}

% ---- Internships & Workshops ------------------------------------------------
\cvsec{Internships \& Workshops}
\yr{2024}{Central Bank of Chile --- Visiting Program}
\yr{2024}{Central Bank of Mexico --- Summer Internship Program}
\yr{2022}{Princeton University --- Macro, Money and Finance --- Continuous Time
  Methods Workshop}

% ---- Other Publications -----------------------------------------------------
\cvsec{Other Publications}
\entrytitlelink{https://ri.conicet.gov.ar/handle/11336/178261}{International Integration and Productive Development: New Policy Guidelines}, 2018\\
with Juan Carlos Hallak --- \textit{Bolet\'in Techint} 356 (in Spanish)
\vspace{0.4em}

\entrytitlelink{https://repositorio.udesa.edu.ar/jspui/handle/10908/17721}{Exports, Bank Credit, and Repayment Capacity of SMEs in Argentina (2007--2016)}, 2020\\
Master thesis, Universidad de San Andr\'es

% ---- Language ---------------------------------------------------------------
\cvsec{Language}
English (professional), Spanish (native), Italian (beginner)

\end{document}
